% RESULTADO_RECUPERADO: funtor dimensional integral, eq:cZ-internos,
% eq:longitud-ruta y definición de ell_0; composición con artículo III.
% La prueba reunida y su control algebraico son específicos de esta contribución.
\subsection{Velocidad constitutiva y realización cinemática del reloj}
\label{v:sec:enlace-velocidad-constitutiva}

La propagación radiativa utiliza la sección de velocidad producida por el
transporte angular del mismo estado HMT. La identificación se establece
antes de elegir unidades de representación. Sean $x=A_{\mathrm{rad}}$ e
$y=C^*_{\mathrm{rad}}$ las coordenadas angulares generadas en la
sección~\ref{v:ant:sec:ii-angulos}, con los canales
\eqref{v:ant:eq:ii-canales}, y sea
$\mathsf R$ la involución autoadjunta de las dos hojas, con proyectores
de rango uno $P_+$ y $P_-$. En el dominio $x>|y|$, el transporte y sus
canales son
\begin{equation}
 T=\exp(-xI-y\mathsf R)=q_+P_++q_-P_-,
 \qquad q_+=e^{-x-y},\quad q_-=e^{-x+y},\quad 0<q_\pm<1.
 \label{v:eq:velocidad-transporte}
\end{equation}
La respuesta de los bloques temporales de longitudes $90$ y $120$ es
\begin{equation}
 \mathsf V=(I-T^{90})(I-T^{120})^{-1}
       =r_+P_++r_-P_-,\qquad
 r_\pm=\frac{1-q_\pm^{90}}{1-q_\pm^{120}}.
 \label{v:eq:velocidad-respuesta}
\end{equation}
La construcción parte de los canales heredados; los símbolos $x,y$ y
$q_\pm$ designan sus coordenadas y no parámetros ajustados a una velocidad
objetivo.

\begin{proposition}[Identidad de las realizaciones constitutiva y cinemática]
\label{v:prop:identidad-velocidad}
En una misma coordenatización dimensional, la sección $c_{\mathrm{int}}$ del reloj
y la sección $c_{\mathrm{vac}}$ de la respuesta constitutiva coinciden:
\begin{equation}
 c_{\mathrm{int}}=c_{\mathrm{vac}}
   =\widehat c\,u_C,
 \qquad
 \widehat c=(r_+r_-)^{-1}
            =\exp(-\mathcal E),
 \label{v:eq:identidad-velocidad}
\end{equation}
donde
\begin{equation}
 \mathcal E=
 \log\!\bigl[(1-q_+^{90})(1-q_-^{90})\bigr]
 -\log\!\bigl[(1-q_+^{120})(1-q_-^{120})\bigr].
 \label{v:eq:velocidad-lector-simetrico}
\end{equation}
La igualdad conserva el intercambio de hojas.
\end{proposition}
\begin{proof}
Los factores que aparecen en los logaritmos son positivos. Por tanto,
\[
 \mathcal E=\log(r_+r_-)
            =\log\det\mathsf V,
 \qquad \exp(-\mathcal E)=(\det\mathsf V)^{-1}.
\]
El determinante se calcula aquí en la realización de dos hojas de rango
uno. La definición cinemática $c_{\mathrm{int}}=\exp(-\mathcal E)u_C$
y la definición constitutiva $c_{\mathrm{vac}}=(r_+r_-)^{-1}u_C$
son, en consecuencia, la misma sección. El intercambio $P_+\leftrightarrow
P_-$ permuta $r_+$ y $r_-$ y conserva su producto. Finalmente,
$\det T=e^{-2x}$ corresponde al transporte elemental, mientras que
$\det\mathsf V=r_+r_-$ corresponde a su respuesta temporal; ambos
determinantes tienen funciones distintas en esta composición.
\end{proof}

\subsubsection{Longitud de trayectoria y duración elemental}

El marco dimensional positivo $(u_S,u_C,u_T,u_Q)$ induce la base de
longitud $u_L=u_Cu_T$. Para una trayectoria enriquecida $\gamma$ de
$n=|\gamma|$ transiciones y canal constante, los lectores aditivos son
\begin{equation}
 T_{\mathrm{int}}(\gamma)=n u_T,
 \qquad L_{\mathrm{int}}(\gamma)=n\widehat c\,u_L.
 \label{v:eq:velocidad-longitud-reloj}
\end{equation}
Su aditividad se refiere al conteo y a la realización de cada transición:
las rutas enriquecidas conservan adicionalmente orientación y memoria.

\begin{corollary}[Normalización de la longitud elemental]
\label{v:cor:longitud-elemental-constitutiva}
Para $n>0$, $L_{\mathrm{int}}(\gamma)/T_{\mathrm{int}}(\gamma)
=c_{\mathrm{vac}}$. En la normalización $t_0=u_T$,
\begin{equation}
 \ell_0=c_{\mathrm{int}}t_0=\widehat c\,u_L.
 \label{v:eq:longitud-elemental-constitutiva}
\end{equation}
En una coordenatización temporal general $t_0=\tau_0u_T$, con $\tau_0>0$,
la expresión correspondiente es $\ell_0=\tau_0\widehat c\,u_L$.
\end{corollary}
\begin{proof}
La división de \eqref{v:eq:velocidad-longitud-reloj} cancela el entero
$n$ y utiliza $u_L/u_T=u_C$. La fórmula elemental se obtiene
multiplicando la sección de velocidad por $t_0$.
\end{proof}

La generalización a un canal dependiente de la arista conserva la fórmula
aditiva $L_{\mathrm{int}}(\gamma)=\sum_{e\in\gamma}\widehat c(e)u_L$.
Su cociente por $nu_T$ es la media de las velocidades de esas aristas.
La realización radiativa homogénea emplea el canal constante declarado
en \eqref{v:eq:velocidad-longitud-reloj}.

\subsubsection{Cambio de coordenatización y normalización adaptada}

\begin{proposition}[Covariancia de la identificación]
\label{v:prop:covariancia-identidad-velocidad}
Sean $u_C'=d u_C$ y $u_T'=\eta u_T$, con $d,\eta>0$.
La misma velocidad y la misma duración elemental tienen coordenadas
\begin{equation}
 \widehat c'=d^{-1}\widehat c,
 \qquad \tau_0'=\eta^{-1}\tau_0,
 \qquad u_L'=d\eta u_L.
 \label{v:eq:velocidad-cambio-carta}
\end{equation}
En particular, si $c_V=\widehat c_Vu_C^V$ y
$u_C^V=d u_C^{III}$, la igualdad de secciones equivale exactamente a
$d\widehat c_V=\widehat c_{III}$.
\end{proposition}
\begin{proof}
Se expresa cada sección en las dos bases. Se obtiene
$\widehat c'u_C'=\widehat c u_C$ y
$\tau_0'u_T'=\tau_0u_T$. Su producto satisface
\[
 \tau_0'\widehat c' u_L'
 =\frac{\tau_0}{\eta}\frac{\widehat c}{d}
       d\eta u_L
 =\tau_0\widehat c u_L=\ell_0.
\]
La última equivalencia resulta de sustituir la relación entre las bases.
\end{proof}

La elección adaptada $d=\widehat c$ da $u_C'=c_{\mathrm{vac}}$ y
$\widehat c'=1$. Esta elección expresa la sección ya construida en una
base adaptada a ella. El coeficiente espectral en la coordenatización interna sigue
siendo $(r_+r_-)^{-1}$; la inversión espectral utiliza esa coordenatización interna.
Por ejemplo, manteniendo las bases de acción y carga, las coordenadas
constitutivas se transforman como
\[
 \widehat\mu'=d\widehat\mu,
 \qquad \widehat\varepsilon'=d\widehat\varepsilon,
 \qquad
 \widehat\mu'\widehat\varepsilon'=(\widehat c')^{-2}.
\]
En la coordenatización adaptada, $\widehat\mu'=r_-/r_+$ y
$\widehat\varepsilon'=r_+/r_-$, pues las coordenadas internas son
$\widehat\mu=r_-^2$ y $\widehat\varepsilon=r_+^2$.

\subsubsection{Momento orientado de Catalán y puente funcional}

La respuesta constitutiva no recibe la constante de Catalán como un
número externo. El cuarto de giro dodecafásico produce la colección de
momentos orientados
\begin{equation}
 \mathcal C_2(s)=\sum_{k\geq0}
 \frac{(-1)^ks^{2k+1}}{(2k+1)^2},\qquad
 \mathcal D=s\frac{\dd}{\dd s},\qquad 0<s<1.
 \label{v:eq:momento-catalan}
\end{equation}
La convergencia absoluta en compactos permite derivar término a término y
da
\begin{equation}
 \mathcal D^2\mathcal C_2(s)=\frac{s}{1+s^2},\qquad
 f(s)=\frac{1+s+s^2}{s(1+s)}\mathcal D^2\mathcal C_2(s)
 =\frac{1-s^3}{1-s^4}.
 \label{v:eq:puente-catalan-constitutivo}
\end{equation}
Con $s_\pm=q_\pm^{30}$, las incidencias $90=3\cdot30$ y
$120=4\cdot30$ producen exactamente
\begin{equation}
 r_\pm=f(s_\pm),\qquad
 \mathcal C_2(1)=G_{\mathrm{Cat}}.
 \label{v:eq:catalan-canales}
\end{equation}
El valor extremo es la representación de la colección; la colección funcional,
y no sólo ese valor, gobierna los dos canales.

El puente es reversible dentro de la clase regular fijada. Dada $f$, la
condición $\mathcal C_2(s)\to0$ cuando $s\downarrow0$ selecciona
unívocamente
\begin{equation}
 \mathcal C_2(s)=\int_0^s\log\frac{s}{t}\,
 \frac{1+t}{1+t+t^2}f(t)\,\dd t.
 \label{v:eq:restitucion-funcional-catalan}
\end{equation}
En efecto, esta integral satisface
$\mathcal D^2\mathcal C_2=s(1+s)f/(1+s+s^2)$; para
\eqref{v:eq:puente-catalan-constitutivo}, el integrando se reduce a
$\log(s/t)/(1+t^2)$. Dos soluciones homogéneas difieren en
$a\log s+b$, y la regularidad nula en el origen fuerza $a=b=0$.

La misma colección dilogarítmica organiza un potencial conjunto. Sea
$\mathcal D_2(z)=\sum_{n\geq1}z^n/n^2$ y
\begin{equation}
 \mathscr F(x,y)=\sum_{\sigma=\pm}
 \left[
 \frac{\mathcal D_2(q_\sigma^{120})}{120}
 -\frac{\mathcal D_2(q_\sigma^{90})}{90}
 \right].
 \label{v:eq:potencial-constitutivo-catalan}
\end{equation}
Como $\dd\mathcal D_2(z)/\dd\log z=-\log(1-z)$ y
$q_\pm=e^{-x\mp y}$,
\begin{equation}
 \partial_x\mathscr F=\log\widehat c,
 \qquad
 \partial_y\mathscr F=\log\widehat Z,
 \qquad \widehat Z=\frac{r_-}{r_+}.
 \label{v:eq:gradiente-constitutivo-catalan}
\end{equation}
Así, velocidad e impedancia son las dos derivadas conjugadas de un mismo
potencial, con constante aditiva fijada por $\mathscr F\to0$ cuando
$q_\pm\to0$.

\subsubsection{Restitución conjunta mediante velocidad y Barbero--Immirzi}

Sea $\psi=f^{-1}$ y definamos
\begin{equation}
 \mathcal H_B(s)=\frac{12s^3}{1-s^9}
 -\frac{s^4}{1-s^{12}}-\frac{(\log s)^2}{20\pi}.
 \label{v:eq:funcional-barbero-espectral}
\end{equation}
La fórmula de \eqref{v:ant:eq:gamma}, junto con
\eqref{v:eq:catalan-canales}, se factoriza exactamente como
\begin{equation}
 \gamma_{\HMT}=\mathcal H_B(\psi(r_-))
 -\mathcal H_B(\psi(r_+)).
 \label{v:eq:factorizacion-barbero-espectral}
\end{equation}

\begin{theorem}[Restitución conjunta de los canales etiquetados]
\label{v:thm:restitucion-c-gamma}
En la colección constitutiva normalizada, el mapa
\begin{equation}
 \Psi:(3/4,1)^2\longrightarrow(1,16/9)\times\mathbb R,
 \qquad
 (r_+,r_-)\longmapsto(\widehat c,\gamma_{\HMT})
 \label{v:eq:mapa-c-gamma}
\end{equation}
es un difeomorfismo. Por tanto, $(\widehat c,\gamma_{\HMT})$ restituye
los dos autovalores etiquetados; los proyectores $P_\pm$ permanecen como
datos de la representación inicial.
\end{theorem}
\begin{proof}
La expresión regular $f(s)=(1+s+s^2)/(1+s+s^2+s^3)$ satisface
\[
 f'(s)=-\frac{s^2(s^2+2s+3)}{(1+s+s^2+s^3)^2}<0,
\]
con $f(0^+)=1$ y $f(1^-)=3/4$. Además,
\begin{equation}
 \mathcal H_B'(s)=
 \frac{36s^2(1+2s^9)}{(1-s^9)^2}
 -\frac{4s^3(1+2s^{12})}{(1-s^{12})^2}
 -\frac{\log s}{10\pi s}>0.
 \label{v:eq:monotonia-hb}
\end{equation}
El cociente del primer término positivo por el segundo antes de restarlo
es
\[
 \frac9s\frac{1+2s^9}{1+2s^{12}}
 \left(\frac{1-s^{12}}{1-s^9}\right)^2>9,
\]
y el último sumando también es positivo. Así,
$F_B=\mathcal H_B\circ\psi$ es estrictamente decreciente, con límites
$F_B(3/4^+)=+\infty$ y $F_B(1^-)=-\infty$.

Fijado $c_0=\widehat c$ y escrito $z=\log\widehat Z$, se tiene
$r_\pm=c_0^{-1/2}e^{\mp z/2}$ y
\[
 |z|<a(c_0):=\min\!\left\{\log c_0,
 \log\frac{16}{9c_0}\right\}.
\]
La función $\Gamma_{c_0}(z)=F_B(r_-)-F_B(r_+)$ cumple
\[
 \Gamma_{c_0}'(z)=\frac12
 \bigl[r_-F_B'(r_-)+r_+F_B'(r_+)\bigr]<0
\]
y recorre $\mathbb R$ entre los dos extremos admisibles. Existe una
solución única para cada $\gamma_{\HMT}$; la derivada no nula y el teorema
de la función implícita proporcionan la inversa suave global.
\end{proof}

Sobre la imagen producida por el generador HMT, la restitución conserva
el orden causal. Recuperados $r_\pm$, se obtienen
$s_\pm=\psi(r_\pm)$ y
\begin{equation}
 A_{\deg}=-\frac3\pi\log(s_+s_-),\qquad
 C^*_{\deg}=\frac3\pi\log\frac{s_-}{s_+},
 \label{v:eq:restitucion-angular-constitutiva}
\end{equation}
de donde, en la orientación positiva ya generada,
$\alpha_{\HMT}=A_{\deg}/1000$ y
$\eta_{\mathrm{ret}}=C^*_{\deg}/2-\alpha_{\HMT}$. Este recorrido
reconstruye representaciones del estado; no convierte
$(\widehat c,\gamma_{\HMT})$ ni valores convencionales en entradas del
generador.

\subsubsection{Transporte al sector radiativo}

La dispersión de la realización radiativa homogénea se escribe
$\omega_{\symbf{k}}=c_{\mathrm{vac}}|\symbf{k}|$.
Con $h=2\pi\hbar$ y $\beta=(k_B\Theta)^{-1}$, su densidad espectral
y su coeficiente de Stefan--Boltzmann admiten las expresiones
\begin{equation}
 u_\omega=
 \frac{\hbar\omega^3}{\pi^2(\widehat c u_C)^3}
 \frac{1}{e^{\beta\hbar\omega}-1},
 \qquad
 \sigma=\frac{\pi^2k_B^4}{60\hbar^3(\widehat c u_C)^2}.
 \label{v:eq:radiacion-velocidad-constitutiva}
\end{equation}
Estas sustituciones transportan exactamente la misma sección de
velocidad. Se conservan las condiciones de la realización radiativa:
sector bosónico libre tridimensional, dos polarizaciones transversales,
potencial químico nulo y temperatura positiva. La identidad
\eqref{v:eq:identidad-velocidad} establece su antecedente constitutivo;
el conteo de modos y el límite termodinámico pertenecen a las pruebas
específicas de dicha realización.

\paragraph{Normalización bajo amplificación.}
Si el refinamiento actúa por $\mathsf V_m=\mathsf V\otimes I_{9^m}$,
la traza normalizada $\tau_m=\operatorname{Tr}/(2\cdot9^m)$ satisface
\[
 \exp[-2\tau_m(\log\mathsf V_m)]
 =\exp[-\log r_+-\log r_-]=\widehat c.
\]
En efecto, cada autovalor aparece $9^m$ veces. La multiplicidad se
cancela en la traza normalizada. El determinante ordinario ampliado es
$(r_+r_-)^{9^m}$: la normalización de dos hojas, o equivalentemente la
traza normalizada, preserva el lector constitutivo durante ese
refinamiento. Esta amplificación de la representación conserva el
transporte enriquecido que la origina.
